% Generated by Listas/make_statements.py from pratica-lecture-08.md, item 13. Do not edit.
\textbf{Source.} Tadelis (2013), \emph{Game Theory: An Introduction}, Exercise 16.3 ``More Limit Pricing'', pp.~335--336 {[}PDF pp.~354--355{]}. Reproduced verbatim. Answer: none in the book.



\begin{quote}

\textbf{More Limit Pricing:} An incumbent firm (player 1) is either a low-cost type \(\theta_1 = \theta_L\) or a high-cost type \(\theta_1 = \theta_H\), each with equal probability. In period \(t = 1\) the incumbent is a monopolist and sets one of two prices \(p_L\) or \(p_H\), and its profits in this period depend on its type and the price it chooses, given by the following table:



{\def\LTcaptype{none} % do not increment counter

\begin{longtable}[]{@{}ccc@{}}

\toprule\noalign{}

Type & Profit from \(p_L\) & Profit from \(p_H\) \\

\midrule\noalign{}

\endhead

\bottomrule\noalign{}

\endlastfoot

\(\theta_L\) & 6 & 8 \\

\(\theta_H\) & 1 & 5 \\

\end{longtable}

}



After observing the period \(t = 1\) price, a potential entrant (player 2), which does not know the incumbent's type but knows the distribution of types, can choose to enter the market (\(E\)) or stay out (\(O\)) in period \(t = 2\). The payoffs of both players in period 2 depend on the entrant's choice and on the incumbent's type and are given by the following table:



{\def\LTcaptype{none} % do not increment counter

\begin{longtable}[]{@{}cccc@{}}

\toprule\noalign{}

Incumbent's type & Entrant's choice & Incumbent's payoff & Entrant's payoff \\

\midrule\noalign{}

\endhead

\bottomrule\noalign{}

\endlastfoot

\(\theta_L\) & \(E\) & 0 & \(-2\) \\

\(\theta_L\) & \(O\) & 8 & 0 \\

\(\theta_H\) & \(E\) & 0 & 1 \\

\(\theta_H\) & \(O\) & 5 & 0 \\

\end{longtable}

}



At the beginning of the game the incumbent discounts profits for period \(t = 2\) using a discount factor \(\delta \le 1\).



\begin{enumerate}

\def\labelenumi{\alph{enumi}.}

\item

  Draw the extensive-form game tree of this game and write down the corresponding matrix.

\item

  For \(\delta = 1\) find a pooling perfect Bayesian equilibrium of the game in which both types of player 1 choose \(p_L\) in period \(t = 1\).

\item

  Find the range of discount factors for which a separating perfect Bayesian equilibrium of the game exists in which type \(\theta_L\) chooses \(p_L\) and type \(\theta_H\) chooses \(p_H\) in period \(t = 1\).

\end{enumerate}

\end{quote}

